% Section 2 -- background. Orientation for a verification reader who may not
% know mathlib's internals. Tight: everything here exists to make Sections 3-6
% readable. No literal numbers: every figure is a macro from numbers.tex.

\section{Background}
\label{sec:background}

\paragraph{mathlib4 and why its dependencies form a graph worth studying.}
Lean~4's mathematical library is a single shared development whose declarations
are each checked against a kernel, so the proof term one elaborates to names
exactly the other declarations it depends on. ``What is used to prove what'' is
recorded mechanically across the whole library, with no citation practice or
bibliographic convention standing between the record and the reasoning.

\paragraph{The published graph.} \citet{li2026network} analyse that structure at
three granularities (declarations, the modules that contain them and their
imports, and namespaces) and release the graph itself. We consume the
declaration layer and do not rebuild it: \graphDeclarations{} declarations and
\graphEdges{} dependency edges, pinned at mathlib commit \mathlibPin{}
(February 2026). Working from the release rather than from our own extraction
is what makes the human side of this paper comparable to theirs.

\paragraph{Most dependency edges were never written by anyone.} The finding that
shapes this work is that only \explicitEdgeCount{} of those edges
(\explicitEdgePct{}\%) are \emph{explicit}; the remaining
\synthesizedEdgePct{}\% are synthesized during elaboration, arriving through
instance resolution, unification and tactic machinery rather than from anything
in the source a contributor typed. A dependency graph read whole is therefore
substantially a record of what the compiler did. Li et al.\ identify the
explicit subgraph as the closer proxy for human-intended dependency structure,
and that subgraph is our object of study throughout. It also settles our
instrument. Source-visible extraction (Section~\ref{sec:extraction}) is the
correct tool for that object rather than a compromise against a fuller one.
What \texttt{explicit} does and does not mean is a separate question, and
Section~\ref{sec:is-explicit} takes it up.

\paragraph{The two corpora, and the difference between them.} Our human
reference is mathlib itself: the \humanProofs{} declarations of the pinned
checkout that carry a located source block. Our machine corpus is a released
set of \machineProofs{} Lean proofs produced by a prover we did not run
\citep{goedelprover}, and its statements come from Lean Workbook
(competition problems). The larger difference is what the two are for. A mathlib declaration is general-purpose
infrastructure written to be reused; a competition proof is a self-contained
puzzle. That difference is a property of the released corpora rather than of
provers, and it is the confound this comparison cannot remove. We state it
here rather than deferring it: Section~\ref{sec:size-controls} narrows it by
matching library area and Section~\ref{sec:limitations}, L4, bounds what the
result can mean in its presence.

\paragraph{Centrality measures technical utility, not mathematical depth.}
Their third finding constrains every structural claim we make. The most heavily
cited declarations in this graph are language infrastructure, and they identify
a two-layer hub structure separating it from mathematical results proper. We
adopt both consequences: no aggregate centrality figure without the layers
separated, and ``more central'' never read as ``less deep'', an inference
\citet{li2026network} contradict in print and
Section~\ref{sec:dag-layer} records us failing to respect.

Li et al.\ pose structural drift, as machine-generated proofs enter the library,
as future work. This paper measures a machine population against their human
baseline.
